% =================================================
% Методическая версия
% =================================================

\worksheettitle
  {Логарифм: в какую степень?}
  {Методическая версия: сценарий, ответы и диагностика понимания}

\begin{teacherbox}[Назначение комплекта]
  Комплект предназначен для первого знакомства с логарифмом. Центральная
  задача урока --- не сообщить новую символику, а помочь ученику увидеть,
  что логарифм является уже знакомым объектом: \textbf{показателем степени}.
  Все задания организованы вокруг перехода
  \[
    a^c=b
    \quad\Longleftrightarrow\quad
    \log_a b=c.
  \]
\end{teacherbox}

\begin{teacherbox}[Методическая граница]
  В этом комплекте условия \(a>0\), \(a\ne1\), \(b>0\) фиксируются и
  применяются к простым записям, но их происхождение не исследуется.
  Доказательное обсуждение ограничений вынесено в отдельный комплект
  «Когда логарифм существует?». Здесь важно не перегрузить первое понятие
  одновременно смыслом, символикой и полной ОДЗ.
\end{teacherbox}

\section{Дидактическая карта}

\begin{center}
  \renewcommand{\arraystretch}{1.35}
  \begin{tabularx}{\textwidth}{@{}>{\bfseries\sffamily}p{0.28\textwidth}X@{}}
    \toprule
    \rowcolor{TableHeader}
    Параметр & Рекомендация \\
    \midrule
    Уровень & 10 класс, первое знакомство с логарифмами \\
    Формат & Один урок 45 минут; расширенный маршрут --- 60--70 минут \\
    Предпосылки & Степени с целыми и дробными показателями, график
      показательной функции, решение простейших показательных уравнений \\
    Основной результат & Ученик переводит запись \(a^c=b\) в
      \(\log_a b=c\) и обратно, вычисляет простые логарифмы как показатели
      степеней \\
    Смысловой результат & Ученик отвечает на вопрос «в какую степень?» до
      выполнения вычислений \\
    Ключевой конфликт & Уравнение \(2^x=5\) имеет единственный корень, но
      его нельзя записать привычным целым или рациональным числом \\
    Основной риск & Формальное запоминание перестановки букв без понимания
      роли основания, аргумента и показателя \\
    Диагностический признак & Ученик проверяет равенство
      \(\log_a b=c\) переходом к \(a^c=b\) \\
    \bottomrule
  \end{tabularx}
\end{center}

\section{Планируемые результаты}

\begin{propertybox}[Предметные результаты]
  К завершению обязательного маршрута ученик:
  \begin{itemize}
    \item объясняет логарифм как показатель степени;
    \item читает и записывает равенство \(\log_a b=c\);
    \item переводит степенную форму в логарифмическую и обратно;
    \item вычисляет логарифмы, если аргумент представим степенью основания;
    \item записывает точный корень уравнения \(a^x=b\) в виде \(\log_a b\).
  \end{itemize}
\end{propertybox}

\begin{teacherbox}[Метапредметный результат]
  Ученик учится распознавать, что новая операция вводится как обратная к
  уже известной. Это важнее количества вычисленных примеров: определение
  становится рабочим инструментом проверки и преобразования записи.
\end{teacherbox}

\section{Маршрут урока}

\begin{center}
  \small
  \renewcommand{\arraystretch}{1.34}
  \begin{tabularx}{\textwidth}{@{}p{0.06\textwidth}p{0.23\textwidth}p{0.14\textwidth}X@{}}
    \toprule
    \rowcolor{TableHeader}
    № & Этап & Время & Действие ученика \\
    \midrule
    1 & Актуализация степеней & 4 мин & Восстанавливает показатели и
      значения степеней: \(5,-2,\frac19,5\) \\
    2 & Проблемная ситуация & 6 мин & Сравнивает уравнения
      \(2^x=8\) и \(2^x=5\) \\
    3 & Графическое подтверждение & 5 мин & Устанавливает существование и
      единственность корня \(2^x=5\) \\
    4 & Введение определения & 8 мин & Формулирует вопрос «в какую степень?»
      и связывает две записи \\
    5 & Двусторонний перевод & 7 мин & Распределяет роли чисел в равенстве \\
    6 & Вычисление логарифмов & 9 мин & Представляет аргумент степенью
      основания \\
    7 & Простейшие уравнения & 4 мин & Переходит к степенной форме \\
    8 & Выходной билет & 2 мин & Объясняет смысл, переводит запись и
      выполняет степенную проверку \\
    \bottomrule
  \end{tabularx}
\end{center}

\begin{teacherbox}[Граница обязательного маршрута]
  За 45 минут целесообразно завершить разделом «Решаем простейшие
  уравнения» и короткой рефлексией. Блоки с анализом ошибок и обобщениями
  \(\log_a1=0\), \(\log_a a=1\) можно использовать как резерв, второй
  фрагмент урока или начало следующего занятия.
\end{teacherbox}

\Needspace{8\baselineskip}
\section{Подробный сценарий}

\stepheading{Этап 1}{Актуализация языка степеней}

Предложите входную диагностику без предварительного повторения. Важно
увидеть, какие виды показателей вызывают затруднения.

\begin{teacherbox}[Ожидаемые ответы]
  \[
    2^5=32,
    \qquad
    10^{-2}=0{,}01,
    \qquad
    3^{-2}=\frac19,
    \qquad
    25^{\frac12}=5.
  \]
\end{teacherbox}

Вопрос классу: «Что общего у найденных чисел?» Ожидаемый ответ:
\textit{все они являются показателями степеней}.

\stepheading{Этап 2}{Создаём потребность в новом обозначении}

Сначала предложите решить \(2^x=8\), затем \(2^x=\frac14\). После этого
запишите \(2^x=5\) и выдержите паузу.

Полезные вопросы:
\begin{itemize}
  \item «Существует ли корень?»
  \item «Между какими целыми числами он расположен?»
  \item «Можно ли назвать его точно без нового обозначения?»
\end{itemize}

Ожидаемая оценка:
\[
  2^2<5<2^3
  \quad\Longrightarrow\quad
  2<x<3.
\]

\begin{teacherbox}[Где важно не подсказывать]
  Не вводите запись \(\log_2 5\) до того, как ученики сформулируют
  нехватку языка: решение есть, оно единственно, но привычного точного
  обозначения пока нет. Тогда логарифм возникает как ответ на реальную
  математическую потребность.
\end{teacherbox}

\stepheading{Этап 3}{Используем график по назначению}

На рисунке ученики должны увидеть ровно одну точку пересечения графиков
\(y=2^x\) и \(y=5\). График нужен не для вычисления точного значения, а
для иллюстрации двух фактов:
\begin{enumerate}
  \item корень существует;
  \item корень единственный.
\end{enumerate}

Отдельно проговорите обоснование единственности: функция \(y=2^x\)
строго возрастает, поэтому каждое своё значение, в частности 5, она
принимает не более одного раза. Рисунок иллюстрирует этот вывод, но не
заменяет рассуждение о монотонности.

Фиксация:
\[
  x_0=\log_2 5.
\]

\begin{warningbox}[Методический риск]
  Если сразу сообщить десятичное приближение \(x\approx2{,}32\), часть
  учеников решит, что логарифм --- это способ приближённого вычисления.
  На первом уроке важнее точная символическая запись и её смысл.
\end{warningbox}

\stepheading{Этап 4}{Формулируем определение}

Сначала попросите закончить фразу:

\begin{center}
  «\(\log_2 5\) --- это такой \ldots, что \(2^{\ldots}=5\)».
\end{center}

Затем замените конкретные числа буквами:
\[
  \log_a b=c
  \quad\Longleftrightarrow\quad
  a^c=b.
\]

Рекомендуемая последовательность вопросов:
\begin{enumerate}
  \item Какое число возводим в степень? --- \(a\).
  \item Какое число хотим получить? --- \(b\).
  \item Что ищем? --- показатель \(c\).
\end{enumerate}

\begin{teacherbox}[Языковая опора]
  На первых заданиях требуйте проговаривать полную фразу:
  «\(\log_3 81\) --- это показатель степени, в которую нужно возвести 3,
  чтобы получить 81». Сокращение речи допустимо только после появления
  устойчивого смысла.
\end{teacherbox}

\stepheading{Этап 5}{Отрабатываем двусторонний перевод}

Не предлагайте мнемоническую перестановку букв как основной метод. Ученик
должен каждый раз распределить три роли:

\begin{center}
  \renewcommand{\arraystretch}{1.3}
  \begin{tabular}{ccc}
    \toprule
    \rowcolor{TableHeader}
    Основание & Показатель & Результат \\
    \midrule
    \(a\) & \(c\) & \(b\) \\
    \bottomrule
  \end{tabular}
\end{center}

После равенства \(5^{-2}=\frac1{25}\) спросите:
«Какое число останется основанием после смены формы записи?»
Ожидаемый ответ: 5.

\stepheading{Этап 6}{Вычисляем как показатели степеней}

Рабочая схема:
\[
  \log_a b
  \longrightarrow
  a^{?}=b
  \longrightarrow
  b=a^c
  \longrightarrow
  \log_a b=c.
\]

Выбирайте примеры, которые различаются типом показателя:
\[
  \log_2 32=5,
  \quad
  \log_5 1=0,
  \quad
  \log_3\frac1{27}=-3,
  \quad
  \log_9 3=\frac12,
  \quad
  \log_{\frac12}8=-3.
\]

\begin{teacherbox}[Диагностический вопрос]
  После каждого ответа просите назвать проверочное степенное равенство.
  Ученик, который говорит только итоговое число, может пользоваться
  подбором без понимания определения.
\end{teacherbox}

После таблицы дайте задание \(\log_{\frac13}27\) без направляющих
пропусков. В полном ответе ученик должен самостоятельно записать вопрос,
получить значение \(-3\) и выполнить проверку
\[
  \left(\frac13\right)^{-3}=27.
\]

\stepheading{Этап 7}{Показываем две роли определения в уравнениях}

Первый тип:
\[
  \log_2 x=6
  \quad\Longleftrightarrow\quad
  x=2^6=64.
\]

Второй тип:
\[
  3^x=10
  \quad\Longleftrightarrow\quad
  x=\log_3 10.
\]

Если классу требуется численная ориентация, покажите результат научного
калькулятора или математической программы:
\[
  \log_3 10\approx2{,}096\approx2{,}10
  \quad\text{(до сотых)}.
\]
Подчеркните, что \(\log_3 10\) --- точное значение, а \(2{,}10\) --- его
приближение. Формулу перехода к новому основанию на этом этапе выводить не
нужно.

Обсудите различие: в первом случае логарифмическая запись переводится в
степенную, во втором --- степенная в логарифмическую.

\stepheading{Этап 8}{Проводим выходную диагностику}

Предложите три коротких задания без обращения к образцам: объяснить смысл
\(\log_2 7\), перевести \(3^{-2}=\frac19\) в логарифмическую форму и
проверить \(\log_4 64=3\) степенным равенством.

Завершите урок вопросом без формулы:

\begin{center}
  «На какой вопрос отвечает логарифм?»
\end{center}

Минимально достаточный ответ:
\textit{«В какую степень нужно возвести основание, чтобы получить данное
число?»}

\section{Ответы к рабочему листу}

\subsection{Входная диагностика}

\[
  5,\qquad -2,\qquad \frac19,\qquad 5.
\]

\subsection{Проблемная ситуация}

\[
  8=2^3\Rightarrow x=3,
  \qquad
  \frac14=2^{-2}\Rightarrow x=-2.
\]

Для \(2^x=5\):
\[
  2^2=4<5<8=2^3,
  \qquad 2<x<3.
\]
Графики имеют одну точку пересечения. Поскольку функция \(y=2^x\) строго
возрастает, значение 5 она принимает только один раз; следовательно,
уравнение имеет единственный корень. По графику нельзя определить его
точное значение.

\subsection{Определение и роли чисел}

Логарифм --- это \textbf{показатель} степени. В записи
\(\log_a b=c\):
\begin{itemize}
  \item \(a\) --- основание степени;
  \item \(b\) --- число, которое нужно получить;
  \item \(c\) --- искомый показатель степени.
\end{itemize}

Главная равносильность:
\[
  \log_a b=c\quad\Longleftrightarrow\quad a^c=b.
\]

В проверке допустимости корректна только запись \(\log_2 8\).
В записи \(\log_{-2}8\) основание отрицательно; в \(\log_2(-8)\)
аргумент отрицателен; в \(\log_1 5\) основание равно 1.

\subsection{Перевод записи}

\[
  5^2=25 \Longleftrightarrow \log_5 25=2,
\]
\[
  10^{-3}=0{,}001 \Longleftrightarrow \lg0{,}001=-3,
\]
\[
  16^{\frac12}=4 \Longleftrightarrow \log_{16}4=\frac12,
\]
\[
  \left(\frac13\right)^{-2}=9
  \Longleftrightarrow
  \log_{\frac13}9=-2.
\]

Обратный перевод:
\[
  4^3=64,
  \qquad
  7^0=1,
  \qquad
  25^{\frac12}=5.
\]

\subsection{Таблица вычислений}

\[
  \log_5 1=0,
  \qquad
  \log_3\frac1{27}=-3,
  \qquad
  \log_9 3=\frac12,
  \qquad
  \log_{\frac12}8=-3.
\]

Самостоятельное задание:
\[
  \left(\frac13\right)^{?}=27,
  \qquad
  27=\left(\frac13\right)^{-3},
  \qquad
  \boxed{\log_{\frac13}27=-3}.
\]
Проверка: \(\left(\frac13\right)^{-3}=27\).

\subsection{Уравнения}

\[
  \log_2 x=6 \Rightarrow x=64,
\]
\[
  \log_5(x+1)=2 \Rightarrow x+1=25 \Rightarrow x=24,
\]
\[
  3^x=10 \Rightarrow x=\log_3 10.
\]

\subsection{Выходной билет}

\begin{enumerate}
  \item \(\log_2 7\) --- показатель степени, в которую нужно возвести 2,
    чтобы получить 7.
  \item \(3^{-2}=\dfrac19\Longleftrightarrow
    \log_3\dfrac19=-2\).
  \item Проверка: \(4^3=64\), следовательно, \(\log_4 64=3\).
\end{enumerate}

\subsection{Анализ ошибок}

Ошибка 1: логарифм является показателем степени, поэтому
\[
  \log_4 64=3,\qquad 4^3=64.
\]

Ошибка 2: основание степени остаётся основанием логарифма:
\[
  2^5=32\Longleftrightarrow\log_2 32=5.
\]

Ошибка 3: значение логарифма может быть отрицательным. Например,
\[
  \log_2\frac14=-2,\qquad 2^{-2}=\frac14.
\]

\subsection{Расширение}

\[
  \log_a1=0,
  \qquad
  \log_a a=1,
  \qquad
  \log_a a^m=m.
\]

Задание 1:
\[
  64=4^3=\left(\frac14\right)^{-3},
  \qquad
  \boxed{\log_{\frac14}64=-3}.
\]

Задание 2:
\[
  \log_3(2x-1)=4
  \Longleftrightarrow
  2x-1=3^4=81
  \Longleftrightarrow
  \boxed{x=41}.
\]

Задание 3:
\[
  \boxed{x=\log_7 20}.
\]
Так как \(7^1=7<20<49=7^2\), то
\[
  1<x<2.
\]

\section{Типичные ошибки и реакция учителя}

\begin{center}
  \small
  \renewcommand{\arraystretch}{1.38}
  \begin{tabularx}{\textwidth}{@{}p{0.26\textwidth}p{0.30\textwidth}X@{}}
    \toprule
    \rowcolor{TableHeader}
    Ошибка & Возможная причина & Реплика или действие учителя \\
    \midrule
    \(\log_4 64=16\) & Логарифм связывают с делением или умножением &
      «Какое степенное равенство проверяет ваш ответ?» \\
    \(\log_{32}2=5\) из \(2^5=32\) & Перепутаны основание и результат &
      Попросить подписать три роли: основание, показатель, результат \\
    Отрицательный ответ отвергается & Степень понимается только как число
      множителей & Вернуться к \(a^{-n}=\frac1{a^n}\) \\
    \(\log_9 3=2\) & Ученик знает \(3^2=9\), но не учитывает направление &
      Записать проверку: должно выполняться \(9^2=3\) \\
    \(3^x=10\) объявляется нерешаемым & Точный ответ отождествляется только
      с рациональным числом & Сравнить обозначения \(\sqrt2\) и
      \(\log_3 10\) \\
    Запись переводится механически & Нет смысловой модели & Требовать устную
      фразу «в какую степень?» перед символическим преобразованием \\
    \bottomrule
  \end{tabularx}
\end{center}

\section{Дифференциация}

\begin{teacherbox}[Базовый маршрут]
  Оставьте только положительные целые, нулевые и отрицательные показатели:
  \(\log_2 8\), \(\log_5 1\), \(\log_3\frac19\). Дробные показатели и
  основание меньше единицы перенесите в домашнюю работу после повторения
  степеней.
\end{teacherbox}

\begin{teacherbox}[Сильная группа]
  Добавьте примеры
  \[
    \log_{\sqrt2}8,
    \qquad
    \log_{\frac1{27}}9,
    \qquad
    \log_{0{,}01}10.
  \]
  Требуйте сначала привести основание и аргумент к одной базе.
\end{teacherbox}

\begin{teacherbox}[Поддержка ученика с пробелами]
  Разрешите пользоваться таблицей степеней чисел 2, 3, 5 и 10. Цель
  первого урока --- не скорость устного счёта, а правильное определение
  показателя и ролей чисел.
\end{teacherbox}

\section{Критерии освоения}

\begin{questionsbox}
  Ученик освоил обязательный материал, если он без подсказки:
  \begin{enumerate}
    \item завершает равносильность
      \(\log_a b=c\Longleftrightarrow a^c=b\);
    \item переводит \(4^{-2}=\frac1{16}\) в логарифмическую форму;
    \item вычисляет \(\log_3\frac1{27}\) и проверяет ответ;
    \item отличает допустимую логарифмическую запись от недопустимой;
    \item самостоятельно оформляет вычисление без направляющих пропусков;
    \item решает \(\log_2 x=5\);
    \item записывает точный корень \(5^x=11\);
    \item словами отвечает на вопрос, что обозначает логарифм.
  \end{enumerate}
\end{questionsbox}

\subsection{Быстрая оценка самостоятельной работы}

\begin{center}
  \renewcommand{\arraystretch}{1.36}
  \begin{tabularx}{\textwidth}{@{}>{\bfseries\sffamily}p{0.23\textwidth}X@{}}
    \toprule
    \rowcolor{TableHeader}
    Уровень & Наблюдаемые признаки \\
    \midrule
    0 --- понятие не сформировано & Ученик использует умножение или деление,
      не может записать проверочную степень \\
    1 --- частичное понимание & Верно считает знакомые примеры, но путает
      основание и аргумент при переводе \\
    2 --- базовое освоение & Верно переводит запись и вычисляет логарифмы с
      целыми показателями \\
    3 --- устойчивое освоение & Работает с отрицательными и дробными
      показателями, объясняет точный ответ в виде логарифма \\
    \bottomrule
  \end{tabularx}
\end{center}

\subsection{Лист наблюдения}

\begin{center}
  \small
  \renewcommand{\arraystretch}{1.7}
  \begin{tabularx}{\textwidth}{@{}p{0.23\textwidth}p{0.15\textwidth}p{0.18\textwidth}p{0.17\textwidth}X@{}}
    \toprule
    \rowcolor{TableHeader}
    Ученик & Вопрос о степени & Перевод записи & Вычисление & Комментарий \\
    \midrule
    & & & & \\
    & & & & \\
    & & & & \\
    & & & & \\
    \bottomrule
  \end{tabularx}
\end{center}

\begin{teacherbox}[Итоговая заметка]
  После урока зафиксируйте, где возник основной разрыв: в степенях, в
  чтении новой записи или в переходе между формами. Это определит, нужен ли
  дополнительный тренинг до перехода к условиям существования логарифма.
\end{teacherbox}
